\section{Newton polygons and a denominator lemma}
\label{sec:newton-denominator}

We first record precisely the local consequence of Newton polygons that
will be used below.  If $L$ is a complete discretely valued field, its
valuation is normalized by $v(L^\times)=\mathbf Z$, and the same letter
$v$ denotes a fixed extension to $\overline{L}$.  We use the standard
Newton polygon theorem over complete discretely valued fields
\cite[Chapter II, \S6, Propositions (6.3) and (6.4)]{Neukirch1999}.

\begin{lemma}[Newton-polygon denominator lemma]
\label{lem:newton-denominator}
Let $P\in L[T]$ be monic with $P(0)\ne0$.  Suppose that the Newton
polygon of $P$ has an edge of horizontal length $\ell$ and slope
\[
 -\frac{h}{d},
 \qquad h\in\mathbf Z,
 \quad d\geq1,
 \quad \gcd(h,d)=1.
\]
Then the slope factor belonging to that edge has degree $\ell$, and
every irreducible factor of the slope factor has degree divisible by
$d$.  In particular, if $\ell=d$, the slope factor is irreducible of
degree $d$ over $L$.
\end{lemma}

\begin{proof}
By the Newton polygon theorem over the complete, hence henselian, field
$L$, the polynomial $P$ has a factor $P_{h/d}\in L[T]$ whose roots,
counted with multiplicity, are exactly the $\ell$ roots $\beta$ of $P$
for which
\[
 v(\beta)=\frac{h}{d}.
\]
Thus $\deg P_{h/d}=\ell$.

Let $G\in L[T]$ be an irreducible factor of $P_{h/d}$ and let $\beta$
be a root of $G$.  Put $M=L(\beta)$ and let $e(M/L)$ be the ramification
index.  With the normalization $v(L^\times)=\mathbf Z$, the value group
of $M$ is
\[
 v(M^\times)=\frac1{e(M/L)}\mathbf Z.
\]
Since $v(\beta)=h/d$ is in lowest terms, it follows that
$d\mid e(M/L)$.  The ramification index divides $[M:L]=\deg G$, so
$d\mid\deg G$.  If $\ell=d$, the degrees of the irreducible factors of
$P_{h/d}$ are positive multiples of $d$ whose sum is $d$; hence there
is exactly one such factor.
\end{proof}

\section{A local quartic factor of the period-4 fiber}
\label{sec:local-quartic-fiber}

Let $\zeta\ne1$ be a root of unity, put $K=\Q(\zeta)$, and write
\[
 \operatorname{ord}(\zeta)=2^a m,
 \qquad m\ \text{odd},
 \qquad \zeta=\xi\eta,
\]
where $\xi$ has $2$-power order and $\eta$ has odd order.  Recall that
\[
 F_\zeta(T)
 =T^6+2T^5+4T^4+6T^3+(\zeta-5)T^2-8T-16.
\]
Fix a prime $\mathfrak p$ of $K$ above $2$, write
$K_{\mathfrak p}$ for the completion, normalize
$v=v_{\mathfrak p}$ by $v(K_{\mathfrak p}^\times)=\mathbf Z$, and put
$e=v(2)$.

When $m>1$, choose the odd-order root of unity $\theta\in K$ used in
Section~\ref{sec:odd-part-base}, so that $\theta^4=\eta$, and put
$t_0=\theta-1$.  The base irreducibility argument already proved that
$t_0$ is a $\mathfrak p$-adic unit and that
\[
 \overline{F_\zeta(T)}=T^2(T-\overline{t_0})^4,
 \qquad
 v\bigl(F_\zeta(t_0)\bigr)=1;
\]
see \eqref{eq:odd-reduction} and Lemma~\ref{lem:t0-valuation}.

\begin{proposition}[Local quartic factor]
\label{prop:local-quartic-factor}
For every root of unity $\zeta\ne1$, the polynomial $F_\zeta$ has an
irreducible factor of degree $4$ over $K_{\mathfrak p}$.
\end{proposition}

\begin{proof}
Suppose first that $m>1$.  Write
\[
 F_\zeta(t_0+Y)=\sum_{j=0}^6 b_jY^j.
\]
Lemma~\ref{lem:t0-valuation} gives $v(b_0)=1$.  Reducing the shifted
polynomial and using \eqref{eq:odd-reduction}, we obtain, in residue
characteristic $2$,
\begin{align*}
 \overline{F_\zeta(t_0+Y)}
 &=(Y+\overline{t_0})^2Y^4 \\
 &=Y^6+\overline{t_0}^{\,2}Y^4.
\end{align*}
Consequently
\[
 v(b_j)\geq1\quad(1\leq j\leq3),
 \qquad v(b_4)=0,
 \qquad v(b_5)\geq1,
 \qquad v(b_6)=0.
\]
Every point $(j,v(b_j))$ lies on or above the line
$y=1-j/4$, and equality holds at $j=0$ and $j=4$.  Hence the Newton
polygon contains the edge
\[
 (0,1)\longrightarrow(4,0),
\]
of length $4$ and slope $-1/4$.  Lemma
\ref{lem:newton-denominator} shows that the corresponding degree-four
slope factor of $F_\zeta(t_0+Y)$ is irreducible over
$K_{\mathfrak p}$.  Translating $Y=T-t_0$ preserves degrees and
irreducibility, and therefore gives the required factor of $F_\zeta$.

Now suppose that $m=1$.  Since $\zeta\ne1$, its order is $2^a$ with
$a\geq1$.  The recalled local cyclotomic facts give
\[
 e=v(2)=2^{a-1}.
\]
For $a=1$ one has $\zeta=-1$ and
$v(\zeta-5)=v(-6)=1$.  For $a\geq2$, the element $\zeta-1$ is a
uniformizer, and
\[
 \zeta-5=(\zeta-1)-4
\]
is a sum of terms of respective valuations $1$ and $2e>1$; hence again
$v(\zeta-5)=1$.  It follows directly from the seven coefficients of
$F_\zeta$, listed from the constant coefficient to the leading
coefficient, that their valuations are
\begin{equation}
 (4e,\,3e,\,1,\,e,\,2e,\,e,\,0).
 \label{eq:F-pure-two-valuations}
\end{equation}

The lower convex hull of these seven points is
\begin{equation}
 (0,4e)\longrightarrow(2,1)\longrightarrow(6,0).
 \label{eq:F-pure-two-newton}
\end{equation}
Indeed, the point $(1,3e)$ lies strictly above the first segment because
\[
 3e>\frac{4e+1}{2},
\]
while the extension of that segment has negative ordinate at every
integer abscissa $j\geq3$, since at $j=3$ its ordinate is
$-2e+3/2<0$.  All the remaining coefficient valuations are
nonnegative.  On the other hand, the points with abscissas $0$ and $1$
lie strictly above the line through $(2,1)$ and $(6,0)$, whose
ordinates there are $3/2$ and $5/4$, respectively, because
$4e>3/2$ and $3e>5/4$.  Finally, the points with abscissas $3,4,5$ lie
strictly above the second segment, whose ordinates there are
$3/4,1/2,1/4$, respectively.  The
two displayed slopes are
\[
 -\frac{4e-1}{2}
 \qquad\text{and}\qquad
 -\frac14,
\]
and are strictly increasing.  These observations also show directly
that no omitted point lies below either segment.  The second edge in
\eqref{eq:F-pure-two-newton} has length $4$ and reduced slope $-1/4$.
Lemma \ref{lem:newton-denominator} therefore makes its slope factor an
irreducible quartic over $K_{\mathfrak p}$.
\end{proof}

\section{A sextic criterion for abelian disjointness}
\label{sec:sextic-abelian-criterion}

Let $k$ be a number field and let
\[
 f(T)=\prod_{i=1}^6(T-r_i)\in k[T]
\]
be an irreducible separable sextic.  For each unordered partition
$\{A,A^c\}$ of $\{1,\ldots,6\}$ into two triples, put
\[
 s_A=\prod_{i\in A}r_i+\prod_{j\in A^c}r_j.
\]
There are $\binom63/2=10$ such partitions.  Define the $3+3$ resolvent
by
\begin{equation}
 \mathcal R_f(S)=
 \prod_{\{A,A^c\}}(S-s_A).
 \label{eq:abstract-three-three-resolvent}
\end{equation}
The absolute Galois group of $k$ permutes the ten unordered partitions;
thus $\mathcal R_f(S)\in k[S]$.

Write $k^{\mathrm{ab}}$ for the maximal abelian extension of $k$ in a
fixed algebraic closure.

\begin{theorem}[Abelian-disjointness criterion for sextics]
\label{thm:sextic-abelian-disjointness}
Let $\alpha$ be a root of $f$ and put $E=k(\alpha)$.  Suppose that
\begin{enumerate}[label=\textup{(\roman*)}]
 \item for some finite place $v$ of $k$, the polynomial $f$ has an
       irreducible quartic factor over $k_v$;
 \item the resolvent $\mathcal R_f$ has no root in $k$.
\end{enumerate}
Then
\[
 k\subseteq D\subseteq E,\quad D/k\text{ Galois}
 \quad\Longrightarrow\quad D=k.
\]
In particular,
\[
 E\cap k^{\mathrm{ab}}=k.
\]
Consequently $f$ remains irreducible over $k^{\mathrm{ab}}$.
\end{theorem}

\begin{proof}
Let $D$ be an intermediate field with $D/k$ Galois.  Since $f$ is
irreducible of degree six,
$[E:k]=6$, and the tower $k\subseteq D\subseteq E$ gives
\[
 [D:k]\mid6.
\]
It remains to exclude the three possibilities $[D:k]=2,3,6$.

Suppose first that $[D:k]=2$, and let $\tau$ be the nontrivial element
of $\Gal(D/k)$.  The minimal polynomial $H$ of $\alpha$ over $D$ is
monic of degree $[E:D]=3$ and divides $f$ in $D[T]$.  The polynomials
$H$ and $\tau(H)$ are distinct: otherwise $H$ would have coefficients
fixed by $\tau$, hence would lie in $k[T]$, contradicting the
irreducibility of the degree-six polynomial $f$ over $k$.  Both are
irreducible over $D$, so they are coprime.  Their product is monic of
degree six, divides $f$, and is fixed by $\tau$; therefore
\[
 f=H\,\tau(H).
\]
The roots of $H$ form a triple $\{r_i:i\in A\}$ and those of
$\tau(H)$ form the complementary triple.  Thus
\[
 H(0)=-\prod_{i\in A}r_i,
 \qquad
 \tau(H)(0)=-\prod_{j\in A^c}r_j,
\]
and consequently
\[
 -\operatorname{Tr}_{D/k}\bigl(H(0)\bigr)
 =\prod_{i\in A}r_i+\prod_{j\in A^c}r_j
 =s_A.
\]
This is a root of $\mathcal R_f$ lying in $k$, contrary to hypothesis
\textup{(ii)}.  Hence $[D:k]\ne2$.

Suppose next that $[D:k]=3$.  Let $q\in k_v[T]$ be an irreducible
quartic factor of $f$, let $\beta$ be a root of $q$, and choose a
$k$-embedding
\[
 \sigma:E\hookrightarrow\overline{k_v}
 \qquad\text{with}\qquad
 \sigma(\alpha)=\beta.
\]
Such an embedding exists because $\beta$ is a root of the minimal
polynomial $f$ of $\alpha$.  Put $L=k_v(\beta)$, so $[L:k_v]=4$.
The compositum $k_v\sigma(D)\subset\overline{k_v}$ is canonically
isomorphic to the completion $D_w$ at the place $w$ induced by
$\sigma$, and it is contained in $L$ because $D\subseteq E=k(\alpha)$.

Since $D/k$ is Galois of prime degree three, the local degree
$[D_w:k_v]$ is either $1$ or $3$.  If it is $3$, then the tower
\[
 k_v\subseteq k_v\sigma(D)\subseteq L
\]
would force $3$ to divide $[L:k_v]=4$, which is impossible.  If it is
$1$, then $v$ splits completely in $D$ and $\sigma(D)\subseteq k_v$.
The minimal polynomial of $\alpha$ over $D$ has degree
$[E:D]=2$; applying $\sigma$ to its coefficients gives a quadratic in
$k_v[T]$ having $\beta$ as a root.  This implies
$[k_v(\beta):k_v]\leq2$, again contradicting $[L:k_v]=4$.  Thus
$[D:k]\ne3$.

Finally, suppose that $[D:k]=6$.  The inclusions $D\subseteq E$ and
the equality of degrees give $D=E$.  In particular, $E/k$ is Galois.
With $q$, $\beta$, and $\sigma$ as in the preceding paragraph,
we have
\[
 k_v(\beta)=k_v\sigma(E),
\]
because $E=k(\alpha)$.  The field on the right is a completion of the
Galois extension $E/k$, so its degree over $k_v$ is the order of a
decomposition group and therefore divides $[E:k]=6$.  Its degree is
also $4$, by the irreducibility of $q$.  This contradiction excludes
$[D:k]=6$.

We have proved that $E$ contains no nontrivial subextension Galois over
$k$.  Now $E\cap k^{\mathrm{ab}}$ is finite and Galois over $k$, so it
equals $k$.  It remains to justify the assertion over the infinite
extension $k^{\mathrm{ab}}$.  If $f$ were reducible over
$k^{\mathrm{ab}}$, a nontrivial factorization would involve only
finitely many coefficients.  They would all lie in some finite
intermediate extension
\[
 k\subseteq M\subseteq k^{\mathrm{ab}}.
\]
The extension $M/k$ is finite abelian Galois, and
\[
 E\cap M\subseteq E\cap k^{\mathrm{ab}}=k.
\]
Thus $E\cap M=k$.  Since $M/k$ is Galois, the fields $E$ and $M$ are
linearly disjoint over $k$.  More explicitly, restriction identifies
$\Gal(EM/E)$ with $\Gal(M/E\cap M)=\Gal(M/k)$; hence
$[EM:E]=[M:k]$, and the tower formula gives
\[
 [M(\alpha):M]=[E:k]=6.
\]
The minimal polynomial of $\alpha$ over $M$ therefore still has degree
six and equals $f$, contradicting the assumed factorization over $M$.
Hence $f$ is irreducible over $k^{\mathrm{ab}}$.
\end{proof}
